\documentclass[../../main.tex]{subfiles}

\begin{document}

\section{自然数与有理数}\label{section:ch01-自然数有理数}

我们想定义自然数为 $1,2,3,\cdots$. 然而, 有必要更为精确. 做这件事的一个方便途径是首先引人归纳集的概念.

\begin{definition}[归纳集]\label{definition:ch01-归纳集}
实数集 $E$ 称为是归纳的, 若它包含 $1$, 且若实数 $x$ 属于 $E$, 则数 $x+1$ 也属于 $E$.
\end{definition}

全体实数集 $\mathbf{R}$ 是归纳的. 从不等式 $1>0$ 我们推出集合 $\{x\in\mathbf{R}\mid x\geqslant 0\}$ 和 $\{x\in\mathbf{R}\mid x\geqslant 1\}$ 是归纳的. 自然数集, 记为 $\mathbf{N}$, 定义为 $\mathbf{R}$ 的所有归纳子集的交. 集合 $\mathbf{N}$ 是归纳的. 为看到这一点, 观察到数 $1$ 属于 $\mathbf{N}$, 这是由于 $1$ 属于每个归纳集. 此外, 若数 $k$ 属于 $\mathbf{N}$, 则 $k$ 属于每个归纳集. 因此, $k+1$ 属于每个归纳集, 所以 $k+1$ 属于 $\mathbf{N}$.

\begin{proposition}[数学归纳法原理]\label{proposition:ch01-数学归纳法}
对每个自然数 $n$, 令 $S(n)$ 为某个数学断言. 假定 $S(1)$ 成立. 也假定每当 $k$ 是使得 $S(k)$ 成立的自然数, 则 $S(k+1)$ 也成立. 那么, 对每个自然数 $n$, $S(n)$ 成立.
\end{proposition}

\begin{proof}
定义 $A=\{k\in\mathbf{N}\mid S(k)\text{ 成立}\}$. 假设恰好意味着 $A$ 是一个归纳集. 于是 $\mathbf{N}\subseteq A$. 因此对每个自然数 $n$, $S(n)$ 成立.
\end{proof}

\begin{theorem}\label{theorem:ch01-自然数最小元}
每个非空自然数集有一个最小成员.
\end{theorem}

\begin{proof}
令 $E$ 为自然数的非空集. 由于集合 $\{x\in\mathbf{R}\mid x\geqslant 1\}$ 是归纳的, 自然数有下界 $1$. 因此 $E$ 有下界 $1$. 作为完备性公理的一个推论, $E$ 有下确界, 定义 $c=\inf E$. 由于 $c+1$ 不是 $E$ 的下界, 存在 $m\in E$ 使得 $m<c+1$. 我们宣称 $m$ 是 $E$ 的最小成员. 否则, 存在 $n\in E$ 使得 $n<m$. 由于 $n\in E$, $c\leqslant n$. 于是 $c\leqslant n<m<c+1$, 且因此 $m-n<1$. 因此自然数 $m$ 属于区间 $(n,n+1)$. 归纳的论证表明对每个自然数 $n$, $(n,n+1)\cap\mathbf{N}=\varnothing$（见习题 8）. 这个矛盾证明了 $m$ 是 $E$ 的最小成员.
\end{proof}

\begin{proposition}[Archimedeas 性质]\label{proposition:ch01-Archimedeas}
对于每对正实数 $a$ 和 $b$, 存在自然数 $n$ 使得 $na>b$.
\end{proposition}

\begin{proof}
定义 $c=b/a>0$. 我们用反证法证明. 若定理是错的, 则 $c$ 是自然数的一个上界. 根据完备性公理, 自然数有一个上确界, 定义 $c_{0}=\sup\mathbf{N}$. 则 $c_{0}-1$ 不是自然数的上界. 选取自然数 $n$ 使得 $n>c_{0}-1$. 因此 $n+1>c_{0}$. 但自然数集是归纳的, 因此 $n+1$ 是自然数. 由于 $n+1>c_{0}$, 而 $c_{0}$ 不是自然数集的上界. 这个矛盾完成了证明.
\end{proof}

我们经常重述 $\mathbf{R}$ 的 Archimedeas 性质: 对每个正实数 $\varepsilon$, 存在自然数 $n$ 使得 $1/n<\varepsilon$.\footnote{Archimedeas 明确表示他的希腊同事 Eurathostenes 确定了我们这里归功于 Archimedeas 的性质.}

定义整数集（记为 $\mathbf{Z}$）为由自然数、它们的相反数和数 $0$ 组成的数集. 有理数集, 记为 $\mathbf{Q}$, 定义为整数的商的集合, 即形如 $x=m/n$ 的数 $x$, 其中 $m$ 和 $n$ 是整数且 $n\neq 0$. 若一个实数不是有理的就称它为无理数. 正如我们在习题 4 证明的, 存在唯一的正数 $x$ 使得 $x^{2}=2$, 记之为 $\sqrt{2}$. 这个数不是有理的. 事实上, 假定 $p$ 和 $q$ 是自然数使得 $(p/q)^{2}=2$, 则 $p^{2}=2q^{2}$. 素数分解定理\footnote{该定理断言每个自然数可唯一地表示为素自然数的乘积, 见\cite{BM97}.}告诉我们 $2$ 除 $p^{2}$ 的次数正好是它除 $p$ 的次数的两倍. 因此 $2$ 除 $p^{2}$ 偶数次. 类似地, $2$ 除 $2q^{2}$ 奇数次. 于是 $p^{2}\neq 2q^{2}$, 且因此 $\sqrt{2}$ 是无理的.

\begin{definition}[在 $\mathbf{R}$ 中稠密]\label{definition:ch01-稠密}
实数的集合 $E$ 称为在 $\mathbf{R}$ 中稠密, 若任何两个实数之间有 $E$ 的成员.
\end{definition}

\begin{theorem}\label{theorem:ch01-有理数稠密}
有理数集在 $\mathbf{R}$ 中稠密.
\end{theorem}

\begin{proof}
令 $a$ 和 $b$ 为实数, 满足 $a<b$. 首先假定 $a>0$. 根据 $\mathbf{R}$ 的 Archimedeas 性质, 存在自然数 $q$ 使得 $(1/q)<b-a$. 再一次利用 $\mathbf{R}$ 的 Archimedeas 性质, 自然数集 $S=\{n\in\mathbf{N}\mid n/q\geqslant b\}$ 非空. 根据\cref{theorem:ch01-自然数最小元}, $S$ 具有最小成员 $p$. 观察到 $1/q<b-a<b$, 于是 $p>1$. 因此 $p-1$ 是自然数（见习题 9）, 因而根据 $p$ 的选取的最小性, $(p-1)/q<b$. 我们也有
\begin{align*}
a=b-(b-a)<(p/q)-(1/q)=(p-1)/q.
\end{align*}
因此有理数 $r=(p-1)/q$ 落在 $a$ 与 $b$ 之间. 若 $a<0$, 根据 $\mathbf{R}$ 的 Archimedeas 性质, 存在自然数 $n$ 使得 $n>-a$. 我们从考虑过的第一种情形推出:存在有理数 $r$ 落在 $n+a$ 与 $n+b$ 之间. 因此有理数 $r-n$ 落在 $a$ 与 $b$ 之间.
\end{proof}

\begin{exercise}\label{exercise:ch01-08}
用归纳法证明: 对每个自然数 $n$, 区间 $(n,n+1)$ 不含任何自然数.
\end{exercise}

\begin{exercise}\label{exercise:ch01-09}
用归纳法证明: 若 $n>1$ 是自然数, 则 $n-1$ 也是一个自然数. 接着用归纳法证明: 若 $m$ 和 $n$ 是满足 $n>m$ 的自然数, 则 $n-m$ 是自然数.
\end{exercise}

\begin{exercise}\label{exercise:ch01-10}
证明对任何实数 $r$, 区间 $[r,r+1)$ 中恰有一个整数.
\end{exercise}

\begin{exercise}\label{exercise:ch01-11}
证明任何有上界的非空整数集有一个最大成员.
\end{exercise}

\begin{exercise}\label{exercise:ch01-12}
证明无理数集在 $\mathbf{R}$ 中稠密.
\end{exercise}

\begin{exercise}\label{exercise:ch01-13}
证明每个实数是某个有理数集的上确界, 也是某个无理数集的上确界.
\end{exercise}

\begin{exercise}\label{exercise:ch01-14}
证明: 若 $r>0$, 则对每个自然数 $n$, $(1+r)^{n}\geqslant 1+n\cdot r$.
\end{exercise}

\begin{exercise}\label{exercise:ch01-15}
用归纳法证明对每个自然数 $n$,
\begin{enumerate}[(i)]
\item $\displaystyle\sum_{j=1}^{n}j^{2}=\frac{n(n+1)(2n+1)}{6}$.
\item $1^{3}+2^{3}+\cdots+n^{3}=(1+2+\cdots+n)^{2}$.
\item 若 $r\neq 1$, $\displaystyle 1+r+\cdots+r^{n}=\frac{1-r^{n+1}}{1-r}$.
\end{enumerate}
\end{exercise}

\end{document}